\chapter{Estado del arte}\label{estado_del_arte}

En esta sección se describen herramientas, aplicaciones y plataformas existentes relacionadas
con el producto propuesto. Además, se exponen investigaciones y estudios recientes que abordan
el análisis del movimiento deportivo mediante inteligencia artificial, visión por computadora
y técnicas de evaluación biomecánica.

\section{Competencia de mercado}

\subsection{VueMotion}

VueMotion \parencite{VueMotionSitio} es una plataforma comercial para análisis del movimiento
deportivo basada en inteligencia artificial y visión por computadora. Se orienta a la
evaluación biomecánica de atletas en contextos reales de entrenamiento, permitiendo analizar
acciones como aceleraciones, carrera, desaceleraciones, cambios de dirección y saltos. La
herramienta se presenta como una solución aplicable a distintos deportes, entre ellos el
fútbol, y busca acercar mediciones biomecánicas que tradicionalmente requerían laboratorios,
sensores o equipamiento especializado.

VueMotion está apuntado a un mercado compuesto por entrenadores, atletas, preparadores
físicos, profesionales médicos, fisioterapeutas, academias deportivas, universidades, clubes
y equipos profesionales. La plataforma comunica sus planes y casos de uso, orientados tanto
a entrenadores individuales como a organizaciones deportivas de mayor escala. Por lo tanto,
su posicionamiento no se limita a usuarios individuales, sino que se enfoca principalmente en
instituciones.

En cuanto a su funcionamiento, VueMotion combina una aplicación móvil con una plataforma web.
La captura puede realizarse mediante dispositivos iOS, y la plataforma permite procesar los
videos para obtener reportes biomecánicos, videos aumentados, kinogramas y datos exportables.
Según la información provista por la empresa, el equipamiento necesario para realizar las
evaluaciones se reduce a un smartphone, un trípode y conos de referencia, sin requerir
sensores corporales ni sistemas de captura de movimiento de laboratorio. No obstante, la
disponibilidad de su aplicación móvil en iOS representa una desventaja para el mercado
argentino, dado que Android concentra el 81,17\% del mercado móvil nacional, frente al
18,79\% correspondiente a iOS \parencite{VueMotionSitio,StatCounterMobileOSArgentina}.

En cuanto al punto de vista técnico, el sitio presenta una similitud importante con el
presente proyecto, ya que ambos parten del uso de video para analizar el movimiento humano
mediante inteligencia artificial y visión por computadora. Esta cercanía se refuerza con la
validación académica realizada por \textcite{DennisonEtAl2026}, donde VueMotion fue
evaluado como un sistema comercial de captura de movimiento sin marcadores basado en
inteligencia artificial para medir variables espaciotemporales del sprint a partir de video
grabado en dispositivos móviles.

No obstante, a partir de la documentación pública disponible, no se identifica con claridad
que VueMotion esté orientado específicamente a la detección automática de patrones de riesgo
de lesión del tren inferior en fútbol. Si bien la plataforma menciona aplicaciones vinculadas
al rendimiento deportivo, la biomecánica, la reducción del riesgo de lesión y el retorno al
juego, no se detallan públicamente reglas, criterios o modelos específicos para detectar
patrones como apoyos inestables, alteraciones en la alineación cadera-rodilla-tobillo o
mecanismos de lesión concretos propios del fútbol o de otros deportes. Su principal
limitación para el presente análisis es que se presenta como una herramienta general de
evaluación biomecánica y rendimiento deportivo, sin evidenciar un foco específico en
patrones preventivos del tren inferior aplicados al fútbol.

\subsection{Zone7}

Zone7 \parencite{Zone7ValidationStudy} es una plataforma comercial basada en inteligencia
artificial orientada a la predicción del riesgo de lesión en deportistas profesionales. A
diferencia de herramientas centradas en el análisis biomecánico por video, su enfoque se
basa en integrar y procesar distintas fuentes de datos disponibles en los clubes, con el
objetivo de generar alertas de riesgo que puedan ser utilizadas por cuerpos técnicos,
preparadores físicos y equipos médicos en la toma de decisiones \parencite{SportsmithZone7}.

El sistema se presenta como una solución que no está condicionada por la tecnología
utilizada por cada institución. Esto significa que puede trabajar con datos provenientes de
diferentes proveedores y dispositivos, principalmente relacionados con carga externa, carga
interna, historial de lesiones, edad del jugador y calendario competitivo. Con estos datos,
esta herramienta estima diariamente la probabilidad de que cada jugador sufra una lesión
dentro de los siguientes siete días, clasificando el riesgo como medio o alto cuando supera
determinados umbrales \parencite{Zone7ValidationStudy}.

Uno de los análisis publicados sobre la herramienta corresponde a un estudio retrospectivo
realizado sobre 11 equipos profesionales de fútbol de Europa y Norteamérica, con datos
pertenecientes desde el 2019 al 2021, donde se cruzaron las predicciones generadas por el
sistema con 423 lesiones reales registradas en jugadores de campo. Según los resultados
publicados, se anticipó un incremento del riesgo entre uno y siete días previos en 306 de
las 423 lesiones analizadas, lo que representa un 72,4\% de los casos
\parencite{Zone7ValidationStudy}.

Dentro del mismo análisis, para nueve de los once equipos, se incluyó un algoritmo orientado
a estimar la zona corporal más probable de lesión en el tren inferior. El sistema indicó las
tres regiones corporales con mayor probabilidad de lesión cada vez que detectaba un aumento
del riesgo. No obstante, los propios autores señalan que esta parte del análisis se encuentra
limitada por el tamaño de muestra disponible por tipo de lesión y que las tendencias
observadas no son estadísticamente significativas \parencite{Zone7ValidationStudy}.

Si bien Zone7 constituye un antecedente relevante por aplicar inteligencia artificial a la
mitigación del riesgo de lesión en fútbol profesional, no se considera una competencia
directa del presente proyecto. Su enfoque principal no se basa en el análisis visual del
movimiento ni en la estimación de postura humana a partir de video, sino en modelos
predictivos construidos sobre datos cargados.

\subsection{Kitman Labs}

Kitman Labs \parencite{KitmanLabsSitio} se orienta a la inteligencia del rendimiento
deportivo mediante la unificación de datos, flujos de trabajo y procesos de toma de
decisión dentro de organizaciones deportivas. Su propuesta incluye soluciones vinculadas
con medicina del rendimiento, optimización del rendimiento, desarrollo de deportistas,
gestión operativa de ligas y análisis personalizado. Su mercado apunta principalmente a
organismos gubernamentales, clubes, equipos, ligas y universidades que necesitan centralizar
información médica, física y deportiva de sus atletas.

Dentro de sus funcionalidades, Kitman Labs incluye un apartado denominado revisión de
lesiones. Esta funcionalidad busca ofrecer una visión clara del tipo, volumen y gravedad de
las lesiones a lo largo del tiempo y durante la temporada. Además, permite analizar factores
como edad, posición y calendario de partidos, con el objetivo de identificar áreas donde
enfocar esfuerzos de prevención, protocolos de retorno y metodología de entrenamiento
\parencite{KitmanLabsAnalisisPersonalizados}.

Al igual que Zone7, su enfoque se apoya principalmente en la integración, organización y
análisis de datos médicos, físicos y contextuales de los deportistas. Por este motivo, si
bien se vincula con la prevención de lesiones y la toma de decisiones en el ámbito
deportivo, no se centra en el análisis visual del movimiento ni en la detección automática
de patrones biomecánicos del tren inferior a partir de video.

\subsection{Catapult}

Catapult \parencite{CatapultInjuryPrevention} se especializa en el monitoreo del rendimiento
físico de los atletas mediante dispositivos wearables, GPS y sensores. En el contexto del
fútbol, sus soluciones permiten registrar y analizar información relacionada con distancia
recorrida, velocidad, aceleraciones, desaceleraciones, cambios de dirección, carga física,
fatiga y recuperación.

Desde el punto de vista preventivo, la empresa propone utilizar estos datos para acompañar
la planificación de cargas de entrenamiento, identificar señales de sobreentrenamiento,
reducir factores asociados al riesgo de lesión y apoyar procesos de retorno al juego. De
esta manera, la información generada puede ser utilizada por entrenadores, preparadores
físicos y cuerpos médicos para tomar decisiones sobre la preparación y disponibilidad de
los jugadores \parencite{CatapultInjuryPrevention}.

Catapult se vincula con la prevención de lesiones a partir del monitoreo físico del jugador.
Su aporte no se basa en analizar visualmente la técnica o la biomecánica del movimiento,
sino en registrar datos mediante GPS, dispositivos wearables y sensores para interpretar la
carga de entrenamiento, la fatiga y la recuperación. A partir de esa información, los
cuerpos técnicos pueden ajustar la planificación, controlar esfuerzos y detectar situaciones
que puedan aumentar el riesgo de lesión \parencite{CatapultInjuryPrevention}.

\section{Herramientas institucionales}

\subsection{NFL Digital Athlete}

Digital Athlete \parencite{NFLDigitalAthlete} es una iniciativa desarrollada por la National
Football League (NFL) junto con Amazon Web Services (AWS) que se presenta como una herramienta
institucional utilizada dentro del ecosistema de la NFL para asistir a sus equipos en la
gestión de la salud y seguridad de los jugadores \parencite{AWSNFL}.

Es una herramienta de predicción de lesiones basada en datos e inteligencia artificial. Su
objetivo es construir una representación integral de la experiencia del jugador, tomando
información proveniente de entrenamientos, prácticas y partidos. Según la NFL, los 32 clubes
de la liga tienen acceso a un portal en el que se presentan datos diarios relacionados con
volumen de entrenamiento, riesgo de lesión, tendencias de la liga y referencias comparativas.
Esta información es utilizada por entrenadores y cuerpos técnicos para diseñar programas
individualizados de prevención, entrenamiento y recuperación.

Desde el punto de vista técnico, el sistema integra múltiples fuentes de datos. AWS describe
que Digital Athlete utiliza información de Next Gen Stats, que son datos de localización,
velocidad y aceleración de los jugadores, videos de partidos y prácticas, sensores,
condiciones climáticas, equipamiento y tipo de jugada. Luego, a partir de estos datos, los
algoritmos de inteligencia artificial y aprendizaje automático ejecutan simulaciones de
partidos y escenarios específicos con el objetivo de identificar situaciones o jugadores con
mayor riesgo de lesión.

\textcite{AWSDigitalAthlete} menciona que esta herramienta se encuentra desarrollando
tecnologías de estimación de pose 3D para analizar cómo los movimientos de los jugadores en
el espacio y el tiempo pueden vincularse con determinadas lesiones. Para ello, utiliza un
conjunto sincronizado de 38 cámaras ubicadas alrededor del estadio, capaces de capturar video
en 5K a 60 cuadros por segundo. Con esta infraestructura, el sistema puede reconstruir
esqueletos virtuales de los jugadores y analizar la posición de articulaciones y extremidades
durante las jugadas.

Esta misma característica marca una de sus principales limitaciones en relación con el
presente análisis, ya que requiere una infraestructura propia de una liga profesional,
compuesta por múltiples cámaras sincronizadas, sensores, datos propietarios de seguimiento y
capacidad de procesamiento en la nube. Por lo tanto, si bien constituye un antecedente
relevante por su uso de inteligencia artificial, datos deportivos, visión por computadora y
prevención de lesiones, no se considera una competencia directa del proyecto propuesto. Su
alcance está orientado al fútbol americano profesional y depende de una infraestructura
avanzada.

\section{Estudios y trabajos realizados}

Dentro de los estudios académicos relevados, se observa una línea de trabajo orientada al
análisis sistemático de videos de lesiones reales ocurridas durante partidos de fútbol. Estos
estudios resultan relevantes para el presente proyecto porque permiten identificar qué
situaciones de juego, mecanismos de lesión y patrones biomecánicos aparecen con mayor
frecuencia en lesiones del tren inferior. A diferencia de las herramientas comerciales
descriptas previamente, estos trabajos no buscan ofrecer una plataforma de uso directo para
clubes o deportistas, sino analizar casos reales de lesión para comprender cómo ocurren y
qué variables podrían considerarse en estrategias de prevención.

Uno de los trabajos más relevantes es el realizado por Della Villa et al.
\parencite{DellaVillaEtAl2020}, quienes analizaron lesiones del ligamento cruzado anterior
en fútbol profesional masculino italiano. El estudio identificó 148 lesiones a lo largo de
diez temporadas, de las cuales 134 pudieron ser analizadas mediante video. Para ello, los
autores utilizaron videos de partidos obtenidos desde plataformas como Wyscout y Panini
Digital, procesaron los fragmentos con herramientas digitales y realizaron la evaluación con
el software Kinovea. A partir de estos videos, tres revisores independientes analizaron el
mecanismo de lesión, el patrón situacional y, cuando la calidad de imagen lo permitía,
variables biomecánicas en el plano frontal y sagital.

Los resultados de dicho trabajo muestran que una gran parte de las lesiones del ligamento
cruzado anterior ocurrieron sin contacto directo sobre la rodilla. El estudio clasificó las
lesiones en mecanismos sin contacto, contacto indirecto y contacto directo, y encontró
cuatro situaciones principales: acciones de presión o entrada, jugador tackleado,
recuperación del equilibrio luego de patear y aterrizaje luego de un salto. Además, la
carga en valgo de rodilla fue identificada como un patrón dominante en la mayoría de los
casos analizados. Este aporte es especialmente importante para el presente proyecto, ya que
permite justificar la observación de variables como alineación de rodilla, apoyo monopodal,
estabilidad del tren inferior y gestos propios del fútbol, como cambios de dirección,
desaceleraciones, entradas y aterrizajes.

En una línea similar, Buckthorpe et al. \parencite{BuckthorpeEtAl2024ACL} analizaron 115
lesiones de ligamento cruzado anterior en fútbol profesional masculino español. Este estudio
tomó como base una metodología comparable a la de Della Villa et al., pero la aplicó sobre
las dos principales categorías del fútbol español. Los videos fueron obtenidos desde
Wyscout, procesados mediante una herramienta en la nube denominada Digital Log y evaluados
con Kinovea. Para el análisis de ángulos y posiciones articulares también se utilizó un
software específico denominado Screen Editor. En este caso, además de los mecanismos de
lesión y los patrones situacionales, los autores incorporaron el análisis de errores
neurocognitivos en las lesiones sin contacto.

Este estudio confirmó varios hallazgos del trabajo anterior. Las lesiones sin contacto y por
contacto indirecto tuvieron una presencia similar, y los patrones más frecuentes volvieron
a estar relacionados con acciones de presión, entradas, jugador tackleado, aterrizajes y
recuperación del equilibrio luego de patear. Desde el punto de vista biomecánico, se
observó un predominio de carga sobre la rodilla, con cadera abducida y valgo dinámico.
Este hallazgo refuerza la importancia de no analizar únicamente una articulación aislada,
sino la relación entre cadera, rodilla, tobillo y apoyo durante acciones específicas del
fútbol. A su vez, la incorporación del componente neurocognitivo muestra que algunas
lesiones ocurren en escenarios donde el jugador debe reaccionar rápidamente ante un
estímulo externo, como un engaño, una marca, una pelota dividida o un cambio inesperado en
la situación de juego.

Otro trabajo relevante para el alcance del presente proyecto es el estudio de Buckthorpe et
al. \parencite{BuckthorpeEtAl2025Ankle}, centrado en esguinces de tobillo en fútbol
masculino de élite. A diferencia de los estudios anteriores, este trabajo no se enfoca en la
rodilla, sino en una lesión frecuente del tobillo, por lo que resulta útil para ampliar el
análisis hacia el tren inferior en su conjunto. Los autores identificaron 166 esguinces de
tobillo y analizaron 140 videos de lesiones ocurridas en partidos oficiales. Al igual que en
los estudios previos, se utilizaron videos obtenidos de plataformas de análisis futbolístico,
procesamiento digital de los clips y evaluación por revisores expertos mediante Kinovea.

Los resultados evidencian, además de lo que ya se viene comentando, que el análisis de video
no solo permite observar la rodilla, sino también la posición del pie, el tipo de apoyo y
la relación entre tobillo y movimiento corporal durante acciones reales de juego.

\textcite{SundbergEtAl2025} realizaron un estudio que sintetizó investigaciones sobre
mecanismos y patrones situacionales de lesión del ligamento cruzado anterior en distintos
deportes, donde se incluyeron varios deportes y más de 5000 situaciones de lesión. Para
evaluar la calidad de los estudios de videoálisis, los autores utilizaron una escala de
evaluación metodológica diseñada por expertos en ciencias del deporte, biomecánica y
medicina deportiva. A partir de la revisión, propusieron una clasificación general de
situaciones de lesión, entre ellas cambios de dirección, aterrizajes, contacto directo sobre
la rodilla y mecanismos inducidos por equipamiento. Estos estudios evidencian que los
mecanismos de lesión no son iguales en todos los deportes, sino que dependen de las demandas
específicas de cada práctica, por lo que no se pueden analizar movimientos genéricos sino
que se deben revisar las acciones propias del fútbol con patrones del tren inferior
asociados a dichas acciones.

Los trabajos de Nassis et al. \parencite{NassisEtAl2023} y Rico-González et al.
\parencite{RicoGonzalezEtAl2023} permiten analizar el uso de machine learning en fútbol
desde una perspectiva más general. Como conclusión puede observarse que el aprendizaje
automático se ha utilizado en áreas como predicción de lesiones, monitoreo de carga,
rendimiento físico, análisis táctico, trayectoria de jugadores y detección de talento.

Estos estudios señalan una limitación en la capacidad predictiva de los modelos de machine
learning aplicados al riesgo de lesión, que hasta la fecha todavía es variable y depende en
gran medida de la cantidad, calidad y continuidad de los datos disponibles. En muchos casos,
los modelos no analizan directamente el movimiento del jugador mediante video, sino que
trabajan sobre datos acumulados del deportista o métricas de carga.

\subsection{Conclusión}

Se puede observar que durante los últimos años aparecieron distintas aproximaciones
tecnológicas y académicas relacionadas con la prevención de lesiones en el fútbol y el
análisis del movimiento deportivo. Algunas soluciones, como Zone7, Kitman Labs y Catapult,
abordan la problemática desde el monitoreo de carga, historial médico, datos físicos, GPS,
sensores y variables contextuales del deportista. Estas herramientas permiten asistir a
cuerpos técnicos y médicos en la toma de decisiones, pero no se centran en el análisis
visual de patrones biomecánicos del tren inferior a partir de video.

Soluciones como VueMotion presentan una mayor cercanía tecnológica con el proyecto
propuesto, ya que utilizan video, inteligencia artificial y análisis biomecánico del
movimiento. Siendo mayores sus similitudes, esta solución se enfoca principalmente en la
evaluación general del rendimiento deportivo y no se evidencia, a partir de la documentación
pública disponible, un foco específico en la detección de patrones de riesgo de lesiones del
tren inferior en fútbol. En el caso de Digital Athlete, se trata de un antecedente relevante
por su uso avanzado de inteligencia artificial, datos deportivos, visión por computadora y
estimación de pose, pero su alcance depende de una infraestructura que requiere una
inversión importante y conjunta de la liga profesional, con múltiples cámaras, sensores y
de la organización de los clubes para compartir y recibir estos datos.

Los estudios académicos relevados permiten reforzar la importancia del análisis de video
para comprender mecanismos reales de lesión en fútbol. En particular, las investigaciones
sobre lesiones del ligamento cruzado anterior y del tobillo muestran que existen patrones
repetidos asociados a situaciones propias del juego, como acciones de presión, cambios de
dirección, desaceleraciones, entradas, aterrizajes, recuperación del equilibrio luego de
patear y apoyos de una sola pierna.

A modo de resumen y teniendo en cuenta el estado del arte expuesto, se identifica una
oportunidad para desarrollar una propuesta orientada al análisis automático de videos de
jugadores de fútbol, enfocada específicamente en el tren inferior y en la detección de
patrones biomecánicos asociados al riesgo de lesión. El proyecto planteado busca
diferenciarse por utilizar videos capturados con dispositivos móviles o cámaras
convencionales, sin requerir sensores corporales ni infraestructura especializada, y por
generar información preventiva sobre rodilla, tobillo y apoyos. Su objetivo, como el de las
herramientas existentes, no es diagnosticar lesiones ni reemplazar al cuerpo médico, sino
aportar una herramienta de apoyo para observar movimientos relevantes y facilitar la
identificación temprana de patrones de riesgo durante acciones propias del fútbol para
poder trabajar en sus causas.
